\section{Representing RDF data in Lean} \label{sec:rdf_in_lean}

This section introduces the Lean declarations used throughout the paper.
Lean 4 provides three constructs for defining types, all used below:
\lean{structure} defines a record type, a fixed set of named fields
grouped into one value (analogous to a class with only data members);
\lean{inductive} defines a type by an exhaustive list of constructors,
each optionally carrying arguments, so a value of the type is built from
exactly one of the listed cases (this is how ShEx's abstract syntax and
semantic rules are represented as data, in Sections~\ref{sec:shex_lean}
and~\ref{sec:semantics_lean}); and \lean{abbrev} introduces a transparent
name for an existing type, used purely for readability. \lean{Finset}~$\alpha$,
from Lean's Mathlib library, is the type of finite sets of elements of
type $\alpha$.

We define an RDF graph as a finite set of RDF triples:

\begin{leanlisting} 
 abbrev Graph := Finset Triple
\end{leanlisting} 

where a \lean{Triple} is a structure:

\begin{leanlisting}
structure Triple where
  subj : Subject
  pred : IRI
  obj  : RDFNode
\end{leanlisting}

such that an \lean{RDFNode} is either an IRI, a literal or a blank node:
\begin{leanlisting}
inductive RDFNode where
  | iri   : IRI → RDFNode
  | lit   : Literal → RDFNode
  | bnode : BNodeId → RDFNode

inductive Subject where
  | iri   : IRI → Subject
  | bnode : BNodeId → Subject
\end{leanlisting}

where IRI, Literal and BNodeId have a simplified definition as:
\begin{leanlisting}
abbrev IRI := String
abbrev BNodeId := String  
inductive Literal where
  | str : String → Literal
  | int : Int → Literal
\end{leanlisting}

A typing is a finite set of pairs of RDF nodes and shape labels:
\begin{leanlisting}
abbrev Assoc  := RDFNode × Label
abbrev Typing := Finset Assoc
\end{leanlisting}

